\chapter{Progettazione}
\label{chap:design}
In questo capitolo vengono illustrate e motivate le principali decisioni architetturali prese durante lo sviluppo del servizio.
Nello specifico, vengono descritte le scelte relative ai componenti che lo costituiscono, con particolare attenzione ai requisiti funzionali e non funzionali individuati nella fase di analisi.
Per alcune delle decisioni prese vengono inoltre discusse le alternative valutate e le motivazioni che hanno portato alla soluzione adottata, così da fornire una visione chiara e completa del percorso progettuale seguito.

\section{Gestione trasparente delle versioni IP}
\label{sec:ipversion-handler}
Il servizio di downstream deve essere in grado di gestire in maniera trasparente le due versioni del protocollo \ac{IP} (IPv4 e IPv6), evitando che la logica applicativa debba occuparsi esplicitamente delle differenze tra i due protocolli.
Le operazioni specifiche di ciascuna versione vengono quindi delegate a una classe dedicata che, sfruttando il pattern \emph{strategy} \cite[pp.~315--323]{gamma1994design}, incapsula le diverse implementazioni necessarie all'invio del messaggio \ac{ICMP} in base alla versione del protocollo \ac{IP} utilizzata.
Rispetto a una soluzione basata su costrutti condizionali (\texttt{if}/\texttt{else} o \texttt{switch}) sparsi nel codice, questa scelta consente di aggiungere in futuro il supporto a nuove versioni del protocollo, o nuove modalità di invio, senza dover modificare il codice client che utilizza tali classi, rispettando così il principio open/closed.

\section{Fan-out parallelo}
\label{sec:parallel-fanout}
L'upstream aggregator ha il compito di eseguire il fan-out\sidenote{Il fan-out si riferisce all'atto di distribuire le informazioni da una singola fonte verso più destinazioni.} verso gli host selezionati.
Poiché il numero di server da interrogare potrebbe aumentare notevolmente nel tempo, non è possibile eseguire le chiamate in sequenza: con un numero elevato di interrogazioni, l'esecuzione verrebbe inevitabilmente interrotta dal timeout imposto. È quindi necessario eseguire le chiamate in parallelo.
L'operazione è di tipo \emph{embarrassingly parallel} \cite{mattson2004patterns}, poiché ogni singola esecuzione è indipendente dalle altre e non richiede alcuna forma di sincronizzazione tra i vari task.
Tuttavia, senza alcuna limitazione, il sistema potrebbe consumare una quantità eccessiva di risorse al crescere del numero di server da interrogare, violando il vincolo introdotto in \Cref{sec:limitazioni}.
Per questo motivo è stato scelto di limitare il numero di unità di esecuzione concorrenti sfruttando il \emph{bounded parallelism pattern}, che attraverso un meccanismo di semafori e monitor limita la concorrenza massima, evitando il degrado della rete e un utilizzo eccessivo delle risorse macchina.

\section{Gestione dei middleware}
\label{sec:middleware-handler}
A seguito di una richiesta \ac{HTTP} potrebbe essere necessario eseguire più funzioni, ciascuna con uno scopo differente. Si vuole poter eseguire tali funzioni in sequenza, con la possibilità di interrompere l'esecuzione della catena in qualsiasi momento, ad esempio nel caso in cui una di esse rilevi una condizione di errore.

Per la gestione dei middleware in sequenza è stato adottato il design pattern \emph{chain of responsibility} \cite[pp.~223--232]{gamma1994design}, che permette di far scorrere la richiesta lungo una catena di ``handler''. Ogni handler, dopo aver processato la richiesta, può decidere se passarla al successivo elemento della catena oppure se interrompere l'esecuzione, restituendo direttamente una risposta.

Alcuni esempi di handler sono la gestione del logging strutturato e il disaster recovery. Sia il servizio di downstream sia quello di upstream utilizzano questa struttura, garantendo un comportamento uniforme e riutilizzabile lungo tutta l'architettura.

\section{Integrazione con FlashStart 2026}
\label{sec:app-integration-design}
Si vuole permettere agli utenti di \emph{FlashStart 2026} di usufruire del servizio di latenza. È dunque necessario aggiungere un nuovo entrypoint all'interno dell'applicazione. Per mantenere una codebase coerente con il resto del sistema, sviluppato dagli altri developer dell'azienda, si è scelto di progettare il nuovo entrypoint seguendo il pattern M.V.C.

L'entrypoint è suddiviso in tre componenti distinti:
\begin{itemize}
    \item il \texttt{LatencyController}, elemento intermediario tra la view e il servizio vero e proprio;
    \item il \texttt{LatencyService}, contenente l'effettiva business logic e il controllo dei permessi;
    \item una serie di repository, che fungono da tramite tra il servizio e le fonti di informazione, quali un database o il servizio di upstream.
\end{itemize}

Al fine di mantenere separata la logica applicativa dalla gestione di autenticazione e autorizzazione, questi ultimi vengono controllati tramite il pattern \emph{decorator} \cite[pp.~175--184]{gamma1994design}, sfruttando le strutture già presenti nella codebase dell'applicazione.

Questa suddivisione, unita alla chiara separazione delle responsabilità tra i vari componenti, garantisce al sistema un'elevata semplicità di manutenzione ed estendibilità nel tempo.